\chapter{Logical Relation for \ThunkLang{}}
\label{chap:logrel}

This chapter describes the development of a logical relation for the
\ThunkLang{} intermediate language of the PureCake compiler. We start with an
analysis of \ThunkLang{}'s semantics and the proof methods currently used for
proving optimisation transformations in it, followed by an overview of logical
relations. Then, we motivate the introduction of a logical relation for
\ThunkLang{} and give our definition along with its properties and applications.

\section{The \ThunkLang{} Intermediate Language}

Each one of PureCake's intermediate languages has its own syntax and semantics.
\ThunkLang{} is the first one, and its structure resembles that of the source
language closely with just two differences: conversion of the calling convention
from call-by-name to call-by-value and introduction of the \lit{force} and
\lit{delay} expressions. We describe its syntax and semantics in depth to help
with the understanding of the logical relation in the next section.

\subsection{Syntax}

\ThunkLang{}'s syntax is described in the following figure:

\begin{figure}[h]
  \begin{bnf}
    $exp$ ::=
    | \lit{var} $x$
    | \lit{Prim} $op$ $\overline{e_n}$
    | \lit{Monad} $mop$ $\overline{e_n}$
    | $e_1 \cdot e_2$
    | $\lambda x.\ e$
    | \lit{letrec} $\overline{x_n = e_n}$ \lit{in} $e$
    | \lit{let} $x = e_1$ \lit{in} $e_2$
    | \lit{if} $e_1$ \lit{then} $e_2$ \lit{else} $e_3$
    | \lit{delay} $e$
    | \lit{force} $e$
    | \lit{value} $v$
    | \lit{tick} $e$
    ;;
  \end{bnf}
  \begin{bnf}
    $v$ ::=
    | \lit{Constructor} C $\overline{v_n}$
    | \lit{Monadic} $mop$ $\overline{e_n}$
    | \lit{Closure} $x$ $e$
    | \lit{Recclosure} $\overline{(x_i, e_i)}$ $x$
    | \lit{Thunk} $e$
    | \lit{Atom} $lit$
    | \lit{DoTick} $v$
    ;;
  \end{bnf}
  \caption{\ThunkLang{} Syntax}
\end{figure}

The $exp$ type describes the language's expressions. We denote the use of
variables by name with \lit{var}. \lit{Prim} and \lit{Monad} apply either a pure
or a monadic operator to a list of expressions. Application and abstraction are
denoted as $e_1 \cdot e_2$ and $\lambda x.\ e$. \lit{let} and \lit{letrec}
denote name binding, \lit{if-then-else} the usual conditional operator, and
\lit{value} embeds a value into an expression. The \lit{delay}, \lit{force} and
\lit{tick} operators are \ThunkLang{}'s special operators. \lit{delay} and
\lit{force} work as described in \autoref{chap:intro} and the tick operator will
be explained in the next section.

The set $v$ describes the language's values. \lit{Constructor} and \lit{Monadic}
construct pure and monadic values respectively, \lit{Closure} and
\lit{Recclosure} represent the closures created by \lit{let} and \lit{letrec},
and \lit{Atom} represents a primitive value. \lit{Thunk} and \lit{DoTick} are
the new values added in \ThunkLang{} when compared with the source language.
\lit{Thunk} wraps an expression in an unevaluated thunk, and \lit{DoTick} will
be explained along with \lit{tick}.

We also observe that $exp$ and $v$ are mutually recursive by means of the
\lit{value} constructor in $exp$. This is due to \ThunkLang{}'s semantics being
substitution based. This means that whenever a function call \hsinl{f} $e$ is
encountered, the resulting value $v$ of evaluating $e$ is substituted into each
appearance of the corresponding formal parameter $p$ of \hsinl{f}, turning every
\lit{var} $p$ expression into \lit{value} $v$. Primary consequence of this is
that when evaluating a \ThunkLang{} expression, encountering a variable name is
always an error, since the only way we can run into a variable expression is if
it hasn't been bound to a value.

\subsection{Semantics}

\ThunkLang{}'s semantics is defined in a big-step functional style
\cite{kahn1987natural}. Big-step semantics is defined by derivation rules for
each construct of the language, and these derivation rules is defined in terms
of a function $|[\cdot|]$ that maps expressions of the language to values.

Usually in big-step semantics, $|[\cdot|]$'s type would be $exp\ \to\ v\ |\
\lit{Fail}$, where \lit{Fail} handles the undefined cases such as encountering a
variable or applying a value that is not a closure. Some derivation rules in
such semantics would be:

\begin{figure}[h]
  \[
    \inference{}{|[\lit{var}\ x|] = \lit{Fail}}\qquad
    \inference
    { |[e_1|] = \lit{true} & |[e_2|] = v}
    { |[\lit{if}\ e_1\ \lit{then}\ e_2\ \lit{else}\ e_3|] = v }
  \]
  \caption{Example big-step evaluation rules}
\end{figure}

Notably, this semantics doesn't capture non-termination. In the
\lit{if-then-else} case, for example, if $|[e_2|]$ doesn't terminate, we have
not defined how $|[\cdot|]$ should behave. Usually this is not problematic, as
translating this semantics to a function in any programming language will imply
termination if any part of the computation terminates.

However, when developing proofs in a formal verification tool such as HOL, every
function defined must be shown to terminate for all possible inputs. Therefore,
writing an evaluation function for any becomes impossible. To mitigate this, we
use a technique called clocking \cite{owens2016functional}. This approach works
by adding a natural number called a clock as an extra argument to $|[\cdot|]_k$,
and whenever the function recurses the clock decreases. The result type is also
augmented with a constructor to represent non-termination, which gets produced
any time the function is called with a zeroed clock. With this transformation,
the evaluation function's type becomes $|[\cdot|]_k\ :\ \mathbb{N}\ \to\  exp\
\to\ v\ |\ \lit{Fail}\ |\ \lit{OOT}$. Proving that $|[\cdot|]_k$ terminates is
now trivial, using the fact that the clock is strictly monotonic decreasing.
With this transformation, the actual semantic derivation rules presented above
for \ThunkLang{} become:

\begin{figure}[h]
  \[
    \inference{}{|[\lit{var}\ x|]_k = \lit{Fail}}\qquad
    \inference
    { |[e_1|]_{k-1} = \lit{true} & |[e_2|]_{k-1} = v }
    { |[\lit{if}\ e_1\ \lit{then}\ e_2\ \lit{else}\ e_3|]_k = v }
  \]
  \[
    \inference
    {}
    { |[\lit{if}\ e_1\ \lit{then}\ e_2\ \lit{else}\ e_3|]_0 = \lit{OOT} }
  \]
  \caption{Example \ThunkLang{} evaluation rules}
\end{figure}

In addition to the clocked function $|[\cdot|]_k$, a regular evaluation function
$|[\cdot|]$ is also usually defined as

\begin{figure}[h]
  \[
    \inference{\nexists k.\ |[e|]_k \neq \lit{OOT}}
    {|[e|] = \lit{OOT}}\qquad
    \inference
    { \exists k.\ |[e|]_k \neq \lit{OOT} & |[e|]_k = v }
    { |[e|] = v }
  \]
  \caption{Top level \ThunkLang{} evaluation function}
\end{figure}


We now have enough context to explain the meaning of the \lit{tick} expression
and \lit{DoTick} value. Similarly to how the \lit{delay} expression wraps its
contained expression in a thunk value for \lit{force} to unwrap, the \lit{tick}
expression wraps its expression in a \lit{DoTick} value. When \lit{force}
encounters a \lit{DoTick} value, it recurses with the underlying value wrapped
in a $\lit{force} \circ \lit{value}$, essentially repeating the force opeation
with one less clock and the unwrapped \lit{DoTick} value.

\begin{figure}[h]
  \[
    \inference
    {|[e|]_k = v}
    { |[\lit{tick}\ e|]_k = \lit{DoTick}\ v }\qquad
    \inference{}{|[\lit{delay}\ e|] = \lit{Thunk}\ e}
  \]
  \[
    \inference{}{|[\lit{force}\ e|]_0 = \lit{OOT}}\qquad
    \inference
    {|[x|]_k = \lit{Thunk}\ e& |[e|]_{k-1} = u}
    {|[\lit{force}\ x|]_k = u}
  \]
  \[
    \inference
    {|[x|]_k = \lit{DoTick}\ v
     & |[\lit{force}\ (\lit{value}\ v)|]_{k-1} = u}
    {|[\lit{force}\ x|]_k = u}
  \]
  \caption{Derivation rules of \ThunkLang{}'s special operators}
\end{figure}

\section{Step-Indexed Logical Relations}

Step-indexed logical relations \cite{dreyer2011logical} are a useful tool for
proving properties of programs. When working with untyped clocked big-step
semantics, proving basic properties about programs becomes repetitive and
error-prone, with reasoning about clocks taking up the majority of the proofs.
Logical relations are already used in CakeML to make reasoning about programs
easier.

A (binary) logical relation relates two constructs that logically match each
other. For example, when considering the results of \ThunkLang{}'s evaluation
function above, we could say that a \lit{Fail} result is related with a
\lit{Fail} result, but not with an \lit{OOT} result. This, of course should
extend to all the constructs of the language. When evaluating two expressions
successfully with results $v_1$ and $v_2$, are the results related or not? It
becomes evident that we need to construct a similar relation for values and for
expressions. After defining the relations for every construct of the language,
we have to prove various compatibility theorems about the relation. These are
the theorems that proofs using the relation depend on to reduce the amount of
work they have to manually do. These compatibility theorems prove that the
relation is closed under the constructs of the language. For example, if $R$ is
a logical relation and we have $R(x_1, x_2)$ and $R(f_1, f_2)$, then $R(f_1
\cdot x_1, f_2 \cdot x_2)$ should hold.

Currently, proofs concerning \ThunkLang{} in PureCake are carried out using
syntax-based relations that require explicit handling of all expression types
and manually matching up the clocks by inserting \lit{tick} expressions. This
adds a lot of boilerplate code obfuscating the essence of the proof, and
requires crafting a new relation every time which comes at the cost of time,
complexity, and duplication. Although logical relations already exist in CakeML,
the user of the relation still needs to match up the clocks when carrying out
proofs. For example, consider the two programs \lit{let $x$ = $e_1$ in $e_2$}
and $e_2[e_1/x]$. These are obviously equivalent, but the first program takes an
extra step compared to the second one. Because the relation is step-indexed,
meaning that programs are not only related to each other but also to a mutual
clock, these two programs are not directly related. This is the reason
intermediate languages in CakeML and PureCake contain \lit{tick} expressions.
The two programs can be show to be equivalent by relating \lit{let $x$ = $e_1$
in $e_2$} and $\lit{tick} (e_2[e_1/x])$. Our goal is to develop a logical
relation that will completely decouple the proofs from the clocking mechanism.

Previous work on developing a logical relation that decouples clocks has been
done in \cite{paraskevopoulou2021compositional}. The authors describe a logical
relation for CertiCoq's $\lambda$ANF intermediate representation that doesn't
require the clocks to match exactly, but accepts any relation over the clocks
that satisfies a set of properties. We take inspiration from their
implementation and adapt it to \ThunkLang{}.

To adapt the logical relation defined in
\cite{paraskevopoulou2021compositional}, we first need to develop an alternative
semantics for \ThunkLang{}, because the clocking in our language is different
from the clocking of their $\lambda$ANF language. As explained in
\cite{kumar2018clocked}, there are two ways of embedding clocks in evaluation
functions: environment-like clocks and state-like clocks. Environment-like
clocking decrements the clocks on each recursive call and doesn't return them,
while state-like clocking passes the clocks around as state, thus limiting
evaluation in terms of length instead of in terms of depth. For example, when
evaluating an \lit{if-then-else} expression with an environment-like clock $k$,
the recursive calls to expressions $e_1$, $e_2$, and $e_3$ will all be passed
the same clock, $k - 1$, whereas with a state like clock, each $e_i$ would get a
clock $k_i$ such that $\sum k_i \le k$. \ThunkLang{}'s semantics use an
environment-like clock.

We define an alternative, yet equivalent, semantics that uses a state-like clock
for the purpose of being compatible with the logical relation. We also add a
second clock called a trace, which only decrements when a function application
occurs instead of at every step. This is used to track function application
depth when proving theorems using the relation about transformations that add or
remove function applications from the programs. This results in the alternative
semantics $|[\cdot|]_k^t$.

\section{A Logical Relation for \ThunkLang{}}

\begin{figure}[p]
\begin{align*}
  &\textbf{Value Relation}\\
  &\mathcal{V}^k(\lit{Constructor $n_1$ $\overline{v_1}$}, \lit{Constructor $n_2$ $\overline{v_2}$})\{Q\} \equiv\\
  &\qquad n_1 = n_2 \land \mathcal{V}^k(\overline{v_1}, \overline{v_2})\{Q\}\\
  &\mathcal{V}^k(\lit{Monadic $op_1$ $\overline{v_1}$}, \lit{Monadic $op_2$ $\overline{v_2}$})\{Q\} \equiv\\
  &\qquad op_1 = op_2 \land \mathcal{V}^k(\overline{v_1}, \overline{v_2})\{Q\}\\
  &\mathcal{V}^k(\lit{Closure $x_1$ $e_1$}, \lit{Closure $x_2$ $e_2$})\{Q\} \equiv\\
  &\qquad \forall i < k\ v_1\ v_2.\ \mathcal{V}^i(v_1, v_2)\{Q\} \land \mathcal{E}^i(e_1[v_1/x_1], e_2[v_2/x_2])\{Q,Q\}\\
  &\mathcal{V}^k(\lit{Recclosure $\overline{f_1}$ $x_1$}, \lit{Recclosure $\overline{f_2}$ $x_2$})\{Q\} \equiv\\
  &\qquad \forall i < k\ e_1.\ e_1 = \overline{f_1}.x_1
  \Rightarrow \exists e_2.\ e_2 = \overline{f_2}.x_2
  \land \mathcal{E}^i(e_1[\overline{f_1}], e_2[\overline{f_2}])\{Q,Q\}\\
  &\mathcal{V}^k(\lit{Thunk $e_1$}, \lit{Thunk $e_2$})\{Q\} \equiv
  \forall i < k.\ \mathcal{E}^i(e_1, e_2)\{Q,Q\}\\
  &\mathcal{V}^k(\lit{Atom $l_1$}, \lit{Atom $l_2$})\{Q\} \equiv
  l_1 = l_2\\
  &\mathcal{V}^k(\lit{DoTick $v_1$}, \lit{DoTick $v_2$})\{Q\} \equiv
  \mathcal{V}^k(v_1, v_2)\{Q\}\\
  &\mathcal{V}^k(v_1, v_2)\{Q\} \equiv \text{False} \qquad\qquad\qquad\qquad \text{Otherwise}\\
  &\textbf{Result Relation}\\
  &\mathcal{R}^k(\lit{Fail}, \lit{Fail})\{Q\} \equiv \text{True}\\
  &\mathcal{R}^k(\lit{OOT}, \lit{OOT})\{Q\} \equiv \text{True}\\
  &\mathcal{R}^k(v_1, v_2)\{Q\} \equiv \mathcal{V}^k(v_1, v_2)\{Q\}\\
  &\mathcal{R}^k(r_1, r_2)\{Q\} \equiv \text{False} \qquad\qquad\qquad\qquad \text{Otherwise}\\
  &\textbf{Expression Relation}\\
  &\mathcal{E}^k(e_1, e_2){Q_L,Q_G} \equiv
  \forall c_1 \le k\ t_1\ r_1.\ |[e_1|]_{c_1}^{t_1} = r_1 \Rightarrow\\
  &\qquad \exists c_2\ t_2\ r_2.\ |[e_2|]_{c_2}^{t_2} = r_2
  \land Q_L(c_1,t_1)(c_2,t_2)
  \land \mathcal{R}^{k-c_1}(r_1,r_2)\{Q_G\}\\
  &\textbf{Variable Relation}\\
  &\mathcal{X}^k(sub_1, sub_2, x, y)\{Q\} \equiv\\
  &\qquad \forall v_1.\ sub_1(x) = v_1 \Rightarrow\\
  &\qquad\qquad\exists v_2.\ sub_2(y) = v_2 \land \mathcal{V}^k(v_1, v_2)\{Q\}\\
  &\textbf{Substitution Relation}\\
  &\mathcal{C}^k(S, sub_1, sub_2)\{Q\} \equiv
  \forall x.\ x \in S \Rightarrow \mathcal{X}^k(sub_1, sub_2, x, x)\{Q\}
\end{align*}
\caption{\ThunkLang{} Logical Relation}
\label{fig:logrel}
\end{figure}

The logical relation, given in whole in \autoref{fig:logrel}, consists of a
value relation $\mathcal{V}$, a result relation $\mathcal{R}$, an expression
relation $\mathcal{E}$, a variable relation $\mathcal{X}$, and a substitution
relation $\mathcal{C}$. Each relation is indexed by a clock k and accepts either
one or two invariants. The expression relation takes two invariats, a local
invariant $Q_L$ and a global invariant $Q_G$, while the others take just one
invariant $Q$.

The value relation $\mathcal{V}$ relates two \ThunkLang{} values. Two
constructor values are related if their constructors are the same and their
argument lists have the same length and are pairwise related. Two closure values
are related if for any pair of values related at a strictly smaller step index
the closure bodies are related at this index after substituting their formal
parameters with the values. Two recursive closure values are related if for any
strictly smaller step index the selected closure bodies are related at that
index after substituting the recursive functions into their bodies. Two thunks
are related if for any strictly smaller step index their enclosed expressions
are related. Finally, two tick expressions are related if for their values are
related.

The result relation $\mathcal{R}$ relates two \ThunkLang{} results. Two
\lit{OOT} results and two \lit{Fail} results are trivially related, and two
value results are related if their values are related.

The substitution relation $\mathcal{C}$ relates two substitutions over a set of
variables $S$, if for any variable in $S$ the substitutions map the variable to
related values, which is expressed by the variable relation $\mathcal{X}$.

\section{Properties}

The invariants the relations range over are binary relations between pairs of
clocks and traces. Using any relation with the equality invariant ($k_1 = k_2
\land t_1 = t_2 \Rightarrow Eq(k_1,t_1)(k_2, t_2)$) is equivalent to using a
standard logical relation. However, more complicated invariants can be used to
relax the restrictions on the clocking. Coming back to the problem of relating
\lit{let $x$ = $e_1$ in $e_2$} and $e_2[e_1/x]$, we can now encode their
relation into the invariants. When working with a syntactic relation as
\ThunkLang{}'s optimisations currently do, a relation between the two
expressions would not work directly. One would have to relate \lit{let $x$ =
$e_1$ in $e_2$} with $\lit{MkTick}(e_2[e_1/x])$ to make the clocks match up,
which complicates both the relation and the proofs needed to show it's sound.
With the logical relation though, one could use a local invariant $Q_L$ defined
as $k_1 = k_2 + 1 \land t_1 = t_2 \Rightarrow Q_L(k_1,t_1),(k_2,t_2)$. Using
this relation, we can directly relate the two expressions since
\[\mathcal{E}^k(\lit{let}\ x = e_1\ \lit{in}\ e_2, e_2[e_1/x])\{Q_L,Eq\}\] This
flexibility is what enables users of the logical relation to decouple their
proofs from the handling of the clocks by coming up with an appropriate
invariant for their transformation.

However, the invariants are subject to some restrictions. We define the
predicate \lit{Compat} $Q_L$ $Q_G$ to mean that the local and global invariants
satisfy a set of compatibility rules. These compatibility rules arise from the
compatibility theorems of the logical relation. Specifically, we need to show
that the relation is closed under \ThunkLang{}'s constructs. For example, we
need to prove a lemma stating that if $\mathcal{E}^k(f_1,f_2)\{Q_L,Q_G\}$ and
$\mathcal{E}^k(x_1,x_2)\{Q_L,Q_G\}$, then $\mathcal{E}^k(f_1 \cdot e_1, f_2
\cdot e_2)\{Q_L,Q_G\}$. For this to be true, we need to assume some relations
about the invariants. Finally, the complete lemma for application compatibility
is:

\newtheorem{theorem}{Theorem}[section]
\newtheorem{corollary}{Corollary}[theorem]
\newtheorem{lemma}[theorem]{Lemma}

\begin{lemma}
  Assume that $Q_L(0,0)(0,0)$ and that
  $Q_L(c_1,t_1)(c_2,t_2)\land Q_L(c_3,t_3)(c_4,t_4)\\
   \Rightarrow Q_L(c_1 + c_3,t_1 + t_3)(c_2 + c_4,t_2 + t_4)$.
  If
  \begin{itemize}
    \item $\mathcal{E}^k(x_1, x_2)\{Q_L,Q_G\}$
    \item $\mathcal{E}^k(f_1, f_2)\{Q_L,Q_G\}$
  \end{itemize}
  then $\mathcal{E}^k(f_1 \cdot x_1, f_2 \cdot x_2)\{Q_L,Q_G\}$
\end{lemma}

We assume that the local invariant holds when both programs time out, and that
it holds for the clocks and traces of the combination of the applications.
Similarly, we need to prove compatibility lemmas for every other language
feature. Each one of these lemmas requires some relation for the invariants, but
the minimum set of relations needed to satisfy all compatibility theorems turns
out to be the following:

\begin{enumerate}

\item $\inference{}{Q_L(0,0)(0,0)}$
  This base condition is required for when both expressions time out.
\item $\inference{Q_L(c_1,t_1)(c_2,t_2)}{Q_L(c_1+1,t_1)(c_2+1,t_2)}$ This
  condition is required for operations that decrease the clock but do not
  perform any applications.
\item $\inference{Q_L(c_1,t_1)(c_2,t_2)}{Q_L(c_1+1,t_1+1)(c_2+1,t_2+1)}$ This
  condition is required for the case where a single application takes place
  along with the decrease in the clock.
\item $\inference{Q_L(c_1,t_1)(c_2,t_2)& Q_L(c_3,t_3)(c_4,t_4)}
                 {Q_L(c_1+c_3,t_1+t_3)(c_2+c_4,t_2+t_4)}$ This last more
  complicated condition is used during the general application case.
\end{enumerate}

\section{Usage}

Currently, each proof in \ThunkLang{} has to define a new syntactical relation
and prove the same lemmas about the relation. Also, an explicit handling of
clocks is required making each proof that adds or removes clock even more
tedious and error-prone. Even worse, \ThunkLang{} has a whole \lit{MkTick}
expression that wraps any other expression and just consumes more clock,
dedicated to aiding in keeping the clocks synchronised across transformations.
Using this logical relation as a framework with appropriate invariants for each
proof, will aid in their development and simplify the language in whole.

We have not yet applied this logical relation to re-write any of the proofs in
\ThunkLang{}. Our work was focused on providing the core of the relation and the
proofs of the compatibility lemmas for the language's constructs. To make the
logical relation practical, a comprehensive summary of the techniques required
for all the proofs must be made, and a library around the relation should be
developed to make the re-writing as smooth as possible, and the results optimal.

